% P15.24: the distinct Stirling-based Durrett Section 3.1 presentation.
\chapter{Stirling's Formula and the de Moivre--Laplace Theorem}
\label{sec:p15-de-moivre-laplace-and-stirling}

This chapter retains the factorial route of Durrett's Section~3.1:
Stirling's formula, local binomial probabilities, lattice sums and
the de Moivre--Laplace limit
\cite[pp.~113--115]{DurrettPTEAuthor2019}.
It also supplies all four exercises in that section, including the
Poisson normal approximation and the two exponential tail rates.
The required Stirling referral is Feller, Volume~I, third edition,
II.9, including its continuation and its later determination of the
constant \cite[pp.~51--54, 179--181]{Feller1968ProbabilityI}.
The logarithmic proof, both remainder bounds and that constant
argument are given in full.

\section{Laws, lattices and the analytic premises}
\label{sec:p15-stirling-premises}

\begin{hbdefinition}{def:p15-stirling-models}{Factorials, sums and normalization conventions}
For an integer \(n\geq1\), \(n!=\prod_{j=1}^n j\), and \(0!=1\).
For integers \(0\leq k\leq n\),
\(\binom nk=n!/(k!(n-k)!)\), with value zero outside that range.
All logarithms in this chapter are natural. For positive sequences,
\(a_n\sim b_n\) means \(a_n/b_n\to1\).

Let \(\varepsilon_i\) be independent fair signs:
\(\mathbb P(\varepsilon_i=1)=\mathbb P(\varepsilon_i=-1)=1/2\).
Write \(S_N=\sum_{i=1}^N\varepsilon_i\), for integers \(N\geq1\).
Its possible values are the lattice
\(\mathcal L_N=\{-N,-N+2,\ldots,N\}\).
Separately, for fixed \(p\in[0,1]\), set \(q=1-p\), and let
\(B_N\) count successes in \(N\) independent Bernoulli trials
of success probability \(p\).
The standard normal density and CDF are
\[
 \phi(x)=(2\pi)^{-1/2}e^{-x^2/2},\qquad
 \Phi(x)=\int_{-\infty}^x\phi(t)\,dt.
\]
Their normalization is the previously proved Gaussian integral,
not a consequence of a normal-limit theorem.
Ceiling and floor denote the least integer above and the greatest
integer below the given real number, respectively.
A supremum of nonnegative approximation errors over an empty
finite index set is assigned the value zero in this chapter.
\end{hbdefinition}

\begin{hblemma}{thm:p15-stirling-log-tools}{proved=annex:p15.stirling-log-tools}{The exact logarithmic estimates and source product limits}
For \(|t|\leq1/2\),
\[
 |\log(1+t)-t|\leq t^2.
\]
For \(0<y<1\),
\[
 \frac1{2y}\log\frac{1+y}{1-y}
       =\sum_{j=0}^\infty\frac{y^{2j}}{2j+1}.
\]
If \(c_n\to0\), \(a_n\to\infty\) and \(a_nc_n\to\ell\in\mathbb R\),
then \((1+c_n)^{a_n}\to e^\ell\), with positive bases eventually.
More generally, for real arrays \(c_{n,j}\), \(1\leq j\leq n\),
\[
 \max_{j\leq n}|c_{n,j}|\to0,\quad
 \sum_{j=1}^nc_{n,j}\to\ell\in\mathbb R,\quad
 \sup_n\sum_{j=1}^n|c_{n,j}|<\infty
 \quad\Longrightarrow\quad
 \prod_{j=1}^n(1+c_{n,j})\to e^\ell.
\]
The last two assertions are source Lemma~3.1.1 and Exercise~3.1.1.
\end{hblemma}

\begin{hblemma}{thm:p15-stirling-tail-premises}{proved=annex:p15.stirling-tail-premises}{The binomial and Gaussian bounds used to determine the constant}
For \(0<p<1\), the binomial mass is
\[
 b(k;N,p)=\mathbb P(B_N=k)=\binom Nk p^kq^{N-k},
 \qquad \mathbb EB_N=Np,\quad\operatorname{Var}(B_N)=Npq.
\]
For integers \(0\leq r\leq N\),
\[
 \mathbb P(B_N\geq r)\leq
 b(r;N,p)\frac{rq}{r-Np}
 \leq\frac{rq}{(r-Np)^2}\quad(r>Np),
\]
and
\[
 \mathbb P(B_N\leq r)\leq\frac{(N-r)p}{(Np-r)^2}
                                             \quad(r<Np).
\]
These are the tail estimates used by Feller's constant argument
\cite[VI.3, pp.~150--151]{Feller1968ProbabilityI}.
For every real \(x>0\), the other required source bounds are
\[
 (x^{-1}-x^{-3})\phi(x)<1-\Phi(x)<x^{-1}\phi(x).
\]
In particular \(1-\Phi(x)\sim\phi(x)/x\) as \(x\to\infty\)
\cite[VII.1, pp.~175, 179]{Feller1968ProbabilityI}.
\end{hblemma}

\section{The complete factorial approximation}
\label{sec:p15-stirling-formula}

\begin{hbtheorem}{thm:p15-stirling-with-remainder}{proved=annex:p15.stirling-with-remainder}{Stirling's formula with the full constant and Robbins bounds}
For every integer \(n\geq1\), there is a real \(r_n\) such that
\[
 n!=\sqrt{2\pi}\,n^{n+1/2}e^{-n}e^{r_n},\qquad
 \frac1{12n+1}<r_n<\frac1{12n}.
\]
Consequently \(n!\sim\sqrt{2\pi n}(n/e)^n\).
The sequence \(r_n\) is strictly decreasing to zero.
The refinement mentioned in Feller's continuation is also valid;
the following explicit stronger remainder is supplied:
\[
 \left|r_n-\frac1{12n}+\frac1{360n^3}\right|
                     \leq\frac6{5n^5}.
\]
The proof first obtains a finite logarithmic constant, then
determines it by Feller's probability-mass argument without
assuming a normal-limit theorem. It also supplies an independent
determination using the already proved cosine-integral central
coefficient.
\index{Stirling's formula}\index{Robbins bounds}
\end{hbtheorem}

\section{The source local formula and the lattice limit}
\label{sec:p15-stirling-local-normal}

\begin{hbtheorem}{thm:p15-stirling-local-rademacher}{proved=annex:p15.stirling-local-rademacher}{Durrett 3.1.2 with an explicit uniform remainder}
For \(n\geq1\) and integers \(-n<k<n\), set \(u=k/n\) and
\[
 H(u)=(1+u)\log(1+u)+(1-u)\log(1-u),\qquad
 R_{n,k}=r_{2n}-r_{n+k}-r_{n-k}.
\]
The exact factorial formula gives
\[
 \mathbb P(S_{2n}=2k)
 =\frac{e^{-nH(u)+R_{n,k}}}{\sqrt{\pi n}\sqrt{1-u^2}}.
\]
For \(0<A<\infty\), \(n\geq\max\{1,2A^2\}\), and
\(x_{n,k}=2k/\sqrt{2n}\) with \(|x_{n,k}|\leq A\),
\[
 \left|\log
 \frac{\mathbb P(S_{2n}=2k)}
 {(\pi n)^{-1/2}e^{-x_{n,k}^2/2}}\right|
 \leq \frac{K_A}{n},
 \qquad K_A=\frac{A^4}{18}+\frac{A^2}{3}+\frac38.
\]
In particular, for integer \(k=k(n)\) with
\(2k/\sqrt{2n}\to x\in\mathbb R\),
\[
 \mathbb P(S_{2n}=2k)\sim(\pi n)^{-1/2}e^{-x^2/2}.
\]
For both parities one has the complete lattice version
\[
 \sup_{\substack{j\in\mathcal L_N\\ |j|\leq A\sqrt N}}
 \left|\frac{\sqrt N}{2}\mathbb P(S_N=j)
              -\phi(j/\sqrt N)\right|\longrightarrow0.
\]
Probabilities at integers of the wrong parity are zero.
At the two extreme lattice points they are exactly \(2^{-N}\);
the interior factorial formula is not applied to \(0!\) through
a nonexistent Stirling remainder \(r_0\).
\end{hbtheorem}

\begin{hbtheorem}{thm:p15-stirling-de-moivre-laplace}{proved=annex:p15.stirling-de-moivre-laplace}{Durrett 3.1.3: the complete probability limit}
For every fixed finite real \(a<b\),
\[
 \mathbb P(a\leq S_N/\sqrt N\leq b)
       \longrightarrow \int_a^b\phi(x)\,dx.
\]
The exact finite probability is
\[
 2^{-N}\sum_{j=L_N}^{U_N}\binom Nj,\qquad
 L_N=\max\!\left\{0,\left\lceil\frac{N+a\sqrt N}{2}\right\rceil\right\},
 \quad
 U_N=\min\!\left\{N,\left\lfloor\frac{N+b\sqrt N}{2}\right\rfloor\right\},
\]
with an empty sum equal to zero.
Replacing either or both interval endpoints by strict inequalities
leaves the limit unchanged. For \(a=b\), the singleton
probability tends to zero. The CDF conclusion is
\[
 \mathbb P(S_N/\sqrt N\leq x)\to\Phi(x)
                          \quad(x\in\mathbb R),
\]
and the interval conclusion also holds for infinite endpoints with
\(-\infty\leq a<b\leq+\infty\).
The source's even-lattice Riemann sums and its odd-index step
\(S_{2n+1}=S_{2n}\pm1\) are both justified in the annex.
\index{de Moivre--Laplace theorem}
\end{hbtheorem}

\section{The Poisson exercise and its probability limit}
\label{sec:p15-stirling-poisson}

\begin{hbdefinition}{def:p15-stirling-poisson-laws}{Poisson laws and their separate sums}
For \(\lambda>0\), the Poisson law of parameter \(\lambda\)
has masses \(e^{-\lambda}\lambda^k/k!\) at integers \(k\geq0\).
For \(\lambda=0\), it is the point mass at zero.
In the following Poisson results, \(X_i\) are independent
Poisson variables of parameter one, and
\(T_n=\sum_{i=1}^nX_i\). This \(T_n\) is separate from the
fair-sign sum \(S_N\).
\index{Poisson distribution}
\end{hbdefinition}

\begin{hblemma}{thm:p15-stirling-poisson-premises}{proved=annex:p15.stirling-poisson-premises}{Normalization, moments and the exact summed law}
Poisson masses sum to one. A Poisson variable of parameter
\(\lambda\geq0\) has mean and variance \(\lambda\).
For the specified independent sequence,
\[
 \mathbb P(T_n=k)=e^{-n}\frac{n^k}{k!}\quad(k\geq0),\qquad
 \mathbb ET_n=\operatorname{Var}(T_n)=n.
\]
All masses at negative integers are zero.
\end{hblemma}

\begin{hbtheorem}{thm:p15-stirling-poisson-normal}{proved=annex:p15.stirling-poisson-normal}{Durrett Exercise 3.1.2 and its complete interval consequence}
If integers \(k=k(n)\geq0\) satisfy
\((k-n)/\sqrt n\to x\in\mathbb R\), then
\[
 \sqrt{2\pi n}\,\mathbb P(T_n=k)\longrightarrow e^{-x^2/2}.
\]
Uniformly over integers with \(|k-n|\leq A\sqrt n\), for
every fixed finite \(A>0\),
\[
 \sqrt n\,\mathbb P(T_n=k)
                  -\phi((k-n)/\sqrt n)\longrightarrow0.
\]
Consequently, for fixed finite \(a<b\),
\[
 \mathbb P(a\leq(T_n-n)/\sqrt n\leq b)
            \longrightarrow\int_a^b\phi(x)\,dx.
\]
The exact sum runs from
\(\max\{0,\lceil n+a\sqrt n\rceil\}\) to
\(\lfloor n+b\sqrt n\rfloor\), and is zero when the upper
index is smaller than the lower.
Strict endpoint variants, singleton limits and infinite endpoint
versions hold as in the preceding fair-sign theorem.
\end{hbtheorem}

\section{The two source large-deviation exercises}
\label{sec:p15-stirling-tail-rates}

\begin{hbdefinition}{def:p15-stirling-rate-functions}{Rate functions and zero probabilities}
For \(-1\leq a\leq1\), define
\[
 \gamma(a)=\tfrac12\{(1+a)\log(1+a)+(1-a)\log(1-a)\},
\]
using \(0\log0=0\).
For \(a>0\), define \(I(a)=a\log a-a+1\).
In logarithms of probabilities, set \(\log0=-\infty\).
These are deterministic rate functions, not probabilities.
\end{hbdefinition}

\begin{hbtheorem}{thm:p15-stirling-rademacher-tail-rate}{proved=annex:p15.stirling-rademacher-tail-rate}{Durrett Exercise 3.1.3 with rounding and endpoints}
For \(0<a<1\),
\[
 \lim_{n\to\infty}\frac1{2n}
       \log\mathbb P(S_{2n}\geq2na)=-\gamma(a).
\]
With \(k_n=\lceil na\rceil\), for all sufficiently large \(n\),
\[
 \mathbb P(S_{2n}=2k_n)
 \leq\mathbb P(S_{2n}\geq2na)
 \leq\frac{\mathbb P(S_{2n}=2k_n)}{1-(1-a)/(1+a)}.
\]
The extended endpoint rates are zero for \(a\leq0\),
\(-\log2\) for \(a=1\), and \(-\infty\) for \(a>1\).
The source assertion is the interior range \(0<a<1\).
\end{hbtheorem}

\begin{hbtheorem}{thm:p15-stirling-poisson-tail-rate}{proved=annex:p15.stirling-poisson-tail-rate}{Durrett Exercise 3.1.4 with the exact geometric comparison}
For \(a>1\),
\[
 \lim_{n\to\infty}\frac1n
       \log\mathbb P(T_n\geq na)
       =-I(a)=a-1-a\log a.
\]
For \(k_n=\lceil na\rceil\),
\[
 \mathbb P(T_n=k_n)
 \leq\mathbb P(T_n\geq na)
 \leq\frac{\mathbb P(T_n=k_n)}{1-1/a}.
\]
For every \(a\leq1\), the corresponding logarithmic rate is zero.
The boundary \(a=1\) is proved separately; its displayed geometric
denominator is not used.
\end{hbtheorem}

\section{An explicit binomial extension and worked boundaries}
\label{sec:p15-stirling-binomial-extension}

\begin{hbtheorem}{thm:p15-stirling-general-binomial}{proved=annex:p15.stirling-general-binomial}{Fixed-parameter binomial normal approximation}
The following extends Durrett's fair-sign formulation.
Fix \(0<p<1\), put \(q=1-p\), and set \(\sigma_N=\sqrt{Npq}\).
Then, for every \(0<A<\infty\),
\[
 \sup_{\substack{k\in\{0,\ldots,N\}\\ |k-Np|\leq A\sigma_N}}
 \left|\sigma_N\mathbb P(B_N=k)
       -\phi((k-Np)/\sigma_N)\right|\longrightarrow0,
\]
and \((B_N-Np)/\sigma_N\) converges in distribution to the
standard normal law. For finite real \(a<b\), the exact interval sum has
indices
\[
 \max\{0,\lceil Np+a\sigma_N\rceil\}
       \leq k\leq\min\{N,\lfloor Np+b\sigma_N\rfloor\}.
\]
At \(p=0\) or \(p=1\), \(B_N=Np\) almost surely and the
variance is zero; the normalizing quotient is not defined.
The theorem fixes \(p\); it makes no claim for arbitrary
sequences \(p=p_N\).
\end{hbtheorem}

\begin{hbtheorem}{thm:p15-stirling-worked-boundaries}{proved=annex:p15.stirling-worked-boundaries}{Numerical approximations and failures of hypotheses}
Let \(A_n=\sqrt{2\pi}\,n^{n+1/2}e^{-n}\), \(n\geq1\).
The relative error \(1-A_n/n!\) is positive and decreases to
zero, although \(n!-A_n\to\infty\).
At \(n=5\), \(A_5\) is approximately \(118.019\), while
\(5!=120\); the improved upper bound \(A_5e^{1/60}\)
is approximately \(120.003\).
The fair-sign rate at \(a=1/2\) is
\(-\gamma(1/2)\approx-0.130812\), and the Poisson rate at
\(a=2\) is \(1-2\log2\approx-0.386294\).

Independence is necessary: if every sign equals the same fair
sign, then \(\mathbb P(|S_N|/\sqrt N\leq1)=0\) for \(N>1\).
Also the bounded scaled-window condition is essential for a
uniform \emph{relative} Gaussian approximation: at \(S_{2n}=2n\),
\[
 \frac{\mathbb P(S_{2n}=2n)}
              {(\pi n)^{-1/2}e^{-n}}
       =\sqrt{\pi n}(e/4)^n\longrightarrow0.
\]
Finally the fixed-\(p\) normal conclusion can fail if \(p_N=1/N\):
the standardized binomial variables then retain an atom of
asymptotic mass \(e^{-1}\) near \(-1\).
\end{hbtheorem}

\begin{proofboundary}{def:p15-stirling-continuation}{The completed source route and its continuation}
This chapter supplies Durrett Section~3.1, including Lemma~3.1.1,
Theorems~3.1.2--3.1.3 and Exercises~3.1.1--3.1.4, through
the specified Stirling route. It includes Feller's II.9
Stirling proof and the complete constant determination it
refers to. The fixed-parameter binomial theorem is an explicit
extension. The separately written weak-convergence theory
supports its distributional interpretation.
Characteristic functions, the general CLT and its later refinements
remain the next distinct source presentations.
\end{proofboundary}
